\section{Declassification Justification and Policy Constraints}

In the decentralized model, strict confidentiality forbids many practically needed releases. Declassification is therefore necessary, but only under explicit authority and explicit target policy \cite{MyersLiskov97}.

\subsection{Authority-Based Declassification Rule}

We use the following principle in server checks:
\[
\frac{\text{authority}(p) \land \text{approved}(T,\theta) \land C_{to} \text{ is justified}}
  {\Gamma \vdash \text{declassify}_{p,T,\theta}(e):(C_{to}, I_{to})}
\]

where authority$(p)$ must cover each owner policy relaxed from source confidentiality to target confidentiality.

This prevents users from widening readers for data they do not control.

\subsection{What Is Declassified in Our System}

We only declassify from clinical-use labels to research-use labels after a trusted transformation pipeline.

Example confidentiality change (same owner):
\[
\{hospital:\{hospital,doctor\}\} \Rightarrow \{hospital:\{hospital,doctor,researcher\}\}
\]

This is permitted only through \texttt{DECLASSIFY\_AND\_RELEASE}, never by ordinary write/relabel operations.

\subsection{Approved Transform Families and Privacy Accounting}

Each transform is approved only if it is both theoretically justified and implemented correctly. We describe each with formal privacy guarantees where applicable:

\textbf{Redaction.} Remove fields that directly identify persons (name, id, address). The transform deletes columns or blanks out values that are quasi-identifiers or directly identifying attributes.

\emph{Why it helps:} Redaction removes the most obvious linkage surface. An attacker who wants to link released records to external datasets must rely on quasi-identifiers (age, gender, zip code) instead of perfect identifiers (name, SSN). This increases the attacker's work and risk of failure.

\emph{Privacy accounting:} Redaction is a preprocessing step that reduces the identifier space but does not add noise. It should be combined with aggregation or $k$-anonymity for formal privacy. Redaction alone is a \emph{heuristic} defense, not a formal guarantee.

\textbf{Aggregation.} Release only group-level summaries (e.g., mean age, count by condition) instead of individual records. The transform groups records and computes aggregate statistics.

\emph{Why it helps:} Individuals are hidden inside population-level values. To reconstruct a single person's record from aggregate queries, an attacker must perform a reconstruction attack, which requires many carefully crafted queries and is limited by the allowed query class.

\emph{Privacy accounting:} The privacy risk depends on the query class. If only mean and count are allowed, an attacker's options are limited. If arbitrary sum queries are allowed, an attacker can potentially reconstruct records using techniques like differential privacy attacks. Policy must restrict query classes to approved patterns.

\textbf{$k$-Anonymization.} Enforce that every combination of quasi-identifier values (age, gender, zip code) appears in at least $k$ records before release. The transform clusters records into equivalence classes and suppresses or generalizes values until each class has at least $k$ members.

\emph{Why it helps:} Bounds the uniqueness risk under quasi-identifiers. If a released record's quasi-identifier matches exactly $k$ people in the population, the attacker cannot identify which individual the record represents (best case) or narrows the candidate set to $k$ people (privacy risk $1/k$).

\emph{Privacy accounting:} For fixed quasi-identifiers $QI$, $k$-anonymity provides a \textbf{worst-case privacy bound}: an attacker with only the quasi-identifiers has at most $1/k$ probability of correctly guessing any individual's sensitive attribute. This is measured as \emph{$k$-anonymity privacy}: $\text{Privacy}(QI, k) = 1/k$.

\emph{Parameter choice:} For hospital data, $k \geq 5$ is typical minimum (ensures at least 5 records per class). For highly sensitive data (e.g., rare diseases), $k \geq 20$ or higher. Policy must enforce $k \geq k_{min}$ where $k_{min}$ is set according to the harm model (how sensitive is disclosure to $1/5$ vs. $1/20$ individuals).

\textbf{Differential Privacy.} Add calibrated random noise to query results such that the output distribution is nearly unchanged whether any single individual's data is included or excluded.

\emph{Why it helps:} Provides a formal, composable privacy guarantee. An attacker observing the query result cannot reliably determine whether the queried database contained any particular individual's record, even with arbitrary background knowledge.

\emph{Privacy accounting:} Under $(\epsilon, \delta)$-differential privacy, for any two datasets $D$ and $D'$ differing by one record, and for any output set $S$:
\[
\mathbb{P}[\text{Mechanism}(D) \in S] \leq e^\epsilon \mathbb{P}[\text{Mechanism}(D') \in S] + \delta
\]

The parameter $\epsilon$ (epsilon) controls privacy: smaller $\epsilon$ means stronger privacy (outputs look more similar regardless of individual's presence). The parameter $\delta$ is the probability of a violation.

\emph{Composition:} Differential privacy composes: if $k$ independent mechanisms each satisfy $(\epsilon_i, \delta_i)$-DP, their composition satisfies $(\sum_i \epsilon_i, k \delta_i)$-DP. This enables privacy budgeting: a hospital can allocate an $\epsilon$ budget per dataset (e.g., $\epsilon \leq 1.0$) and track consumption as queries are answered. Once $\epsilon$ is exhausted, no further queries are answered on that dataset.

\emph{Parameter choice:} The choice of $\epsilon$ reflects the privacy-utility tradeoff. $\epsilon = 1.0$ (often written $\epsilon \leq \log(3)$) provides moderate privacy with acceptable utility. $\epsilon = 0.1$ provides strong privacy with noticeable noise. $\epsilon = 10$ provides weak privacy. Policy must enforce $\epsilon \leq \epsilon_{max}$ where $\epsilon_{max}$ is set to ensure total budget does not exceed total allocated privacy loss (e.g., per-hospital privacy budget of $\epsilon_{total} = 5.0$ with per-study allocations).

\subsection{Policy Parameters (Concrete and Checkable)}

To avoid superficial "approved transform" claims, the policy must define concrete thresholds:
\begin{itemize}
  \item minimum group size ($k \geq k_{min}$),
  \item maximum privacy budget per dataset ($\epsilon \leq \epsilon_{max}$),
  \item allowed query classes (no unrestricted row-level predicates),
  \item logging requirements (who approved, why, and on what dataset version).
\end{itemize}

These constraints are as important as the transform itself. A weak parameter choice can nullify theoretical guarantees.

\subsection{Why This Is Compatible with Myers-Liskov}

Our design follows the paper's declassification philosophy in three ways:
\begin{itemize}
  \item declassification is explicit and local (single API command),
  \item declassification requires authority to relax owner policies,
  \item declassified result receives a new explicit label and is auditable.
\end{itemize}

We do \emph{not} allow hidden, implicit, or user-defined declassification semantics.

\subsection{Threat-Oriented Justification}

We justify declassification against likely attacker strategies:
\begin{itemize}
  \item \textbf{Linkage attack}: combine released fields with external datasets.
  \item \textbf{Differencing attack}: compare near-identical aggregate queries.
  \item \textbf{Repeated query attack}: consume privacy budget over time.
\end{itemize}

Mitigations are policy-coupled:
\begin{itemize}
  \item redaction and generalization reduce linkage surface,
  \item minimum cohort size and query restrictions reduce differencing precision,
  \item cumulative budget tracking and cutoff enforce bounded repeated leakage.
\end{itemize}

\subsection{Hospital Scenario Example: Multi-Step Privacy Pipeline}

Consider patient data with label $L_{\text{clinical}} = (patient:\{hospital,doctor\}) \wedge (hospital:\{hospital\})$. For research access, the hospital applies:

\begin{enumerate}
  \item \textbf{Redaction:} Remove direct identifiers (name, SSN, MRN).
  \item \textbf{Aggregation:} Release group statistics instead of individual records.
  \item \textbf{$k$-Anonymity:} Verify quasi-identifier combinations appear in at least $k = 5$ records.
  \item \textbf{Authority Check:} Administrator confirms researcher authorization and ethics approval.
  \item \textbf{Release and Audit:} Create output with label $L_{\text{research}} = (patient:\{hospital,doctor,researcher\}) \wedge (hospital:\{hospital,researcher\})$.
\end{enumerate}

This pipeline combines redaction (reduces linkage surface), aggregation (hides individuals), $k$-anonymity (formal privacy bound $1/k$), authority verification (only authorized release), and audit (retroactive accountability).

